% !TEX root = IE3_expC_report.tex
\documentclass[lualatex,a4paper,10pt]{jlreq}
% Shared LuaLaTeX preamble.
\usepackage{luatexja-fontspec}
\setmainfont{TeX Gyre Termes}
\setsansfont{TeX Gyre Heros}
\setmonofont{Latin Modern Mono}
\setmainjfont[BoldFont=HaranoAjiGothic-Medium]{HaranoAjiMincho-Regular}
\setsansjfont[BoldFont=HaranoAjiGothic-Bold]{HaranoAjiGothic-Medium}
\usepackage[top=22mm,bottom=22mm,left=23mm,right=23mm]{geometry}
\usepackage{amsmath,amssymb,booktabs,array,longtable,tabularx}
\usepackage{graphicx,xcolor,tikz,circuitikz}
\usetikzlibrary{arrows.meta,positioning,calc}
\usepackage{enumitem,fancyhdr,listings,needspace}
\usepackage[unicode,colorlinks=true,linkcolor=labblue,urlcolor=labteal]{hyperref}
\usepackage{bookmark}
\definecolor{labblue}{RGB}{30,70,112}
\definecolor{labteal}{RGB}{0,100,105}
\definecolor{labgray}{gray}{0.95}
\ModifyHeading{section}{font={\large\sffamily\bfseries\color{labblue}}}
\ModifyHeading{subsection}{font={\normalsize\sffamily\bfseries\color{labteal}}}
\setlength{\parindent}{1\zw}
\setlength{\parskip}{3pt}
\setlength{\emergencystretch}{2em}
\renewcommand{\arraystretch}{1.35}
\setlist{itemsep=3pt,topsep=4pt,leftmargin=2\zw}
\newcolumntype{Y}{>{\raggedright\arraybackslash}X}
\newcolumntype{C}[1]{>{\centering\arraybackslash}p{#1}}
\pagestyle{fancy}
\fancyhf{}
\fancyhead[L]{\small\color{labblue} IE3 後期・電子工学実験}
\fancyhead[R]{\small テーマ\LabCode}
\fancyfoot[C]{\thepage}
\setlength{\headheight}{16pt}
\DeclareOldFontCommand{\rm}{\normalfont\rmfamily}{\mathrm}
\lstset{basicstyle=\small\ttfamily,columns=fullflexible,keepspaces=true,
 breaklines=true,frame=single,backgroundcolor=\color{labgray},
 numbers=left,numberstyle=\tiny,showstringspaces=false,aboveskip=8pt,belowskip=8pt}
\newcommand{\blank}{\rule{22mm}{0.3pt}}
\newcommand{\answerline}{\par\noindent\rule{\linewidth}{0.25pt}\par}
\newcommand{\writing}[1]{\par\Needspace{\dimexpr#1+5mm\relax}\noindent%
 \fcolorbox{labblue!35}{white}{\parbox[c][#1][t]{\dimexpr\linewidth-2\fboxsep-2\fboxrule\relax}{\vspace{2mm}\noindent\strut\hfill\par}}\par}
\newcommand{\hint}[1]{\par\smallskip\noindent{\sffamily\bfseries\color{labteal}記述の視点：}#1\par}
\newcommand{\note}[1]{\par\smallskip\noindent\fcolorbox{labblue!45}{labblue!5}{%
 \parbox{\dimexpr\linewidth-2\fboxsep-2\fboxrule\relax}{#1}}\par\smallskip}
\newcommand{\eqerror}{\varepsilon=\frac{x_{\mathrm{meas}}-x_{\mathrm{th}}}{x_{\mathrm{th}}}\times100\ \%}
\newcommand{\LabSource}{徳山工業高等専門学校 情報電子工学科，
 『2026年度 電子工学実験（3年後期）テキスト』，テーマ\LabCode，
 \expandafter\path\expandafter{\SourcePdf}。}
\newcommand{\reporttitle}{
 \noindent{\small\sffamily\color{labteal}実験報告書・記入ガイド}\par
 \noindent{\LARGE\sffamily\bfseries \LabCode\quad\LabTitle}\par\smallskip
 \noindent 氏名：\blank\quad 班：\blank\quad 実験日：\blank\par
 \note{目的・原理・手順を確認し、実測値と自分の考察を記入する。
 考察のガイドは、検討する事象・使う測定結果・記述の方針を示す。}
}
\newcommand{\prelabtitle}{
 \noindent{\small\sffamily\color{labteal}事前学習・前提知識プリント}\par
 \noindent{\LARGE\sffamily\bfseries \LabCode\quad\LabTitle}\par\smallskip
 \noindent 氏名：\blank\quad 班：\blank\quad 予習日：\blank\par
 \note{用語、回路の動き、計算、測定表への記入の順に読む。
 例題の値は説明用であり、実験結果ではない。}
}
\newcommand{\tocrow}[3]{\hyperref[#1]{#2}& #3 &\pageref*{#1}\\[2mm]}
\newcommand{\documentmap}[1]{
 \section*{目次と概要}
 \noindent 青色の項目をクリックすると該当箇所へ移動する。\par
 {\small\begin{longtable}{@{}p{0.29\linewidth}p{0.61\linewidth}r@{}}
 \toprule {\color{labblue}\bfseries 項目} & {\color{labblue}\bfseries 読むと分かること} & 頁\\\midrule
 \endhead #1\bottomrule
 \end{longtable}}
 \clearpage
}
\newenvironment{terms}{\par\smallskip\begin{longtable}{@{}p{0.23\linewidth}p{0.73\linewidth}@{}}
 \toprule {\color{labteal}\bfseries 用語}&{\color{labteal}\bfseries 意味・回路での役割}\\\midrule\endhead}
 {\bottomrule\end{longtable}}
\newcommand{\term}[2]{\textbf{#1}&#2\\[2mm]}
\newcommand{\guidepoint}[4]{
 \Needspace{32mm}\subsection{#1}
 \noindent{\bfseries\color{labteal}考える事象：}#2\par
 \noindent{\bfseries\color{labteal}参照する結果：}#3\par
 \noindent{\bfseries\color{labteal}記述の方針：}#4\par
}
\newcommand{\reportclosing}[1]{
 \clearpage\section{結論のまとめ方}\label{sec:closing}
 \hint{#1}
 \ConclusionGuide
 \begin{enumerate}
 \item 目的に対応する最も重要な実測結果を、測定条件と数値を添えて述べる。
 \item 原理と一致した点・一致しなかった点を示す。原因は確認済みの事実と推測を分ける。
 \item 実施範囲で言えることと、未確認の条件を一文ずつまとめる。
 \end{enumerate}
 \note{書き出しの型：「○○の条件で△△を測定したところ、□□となった。
 これは◇◇と整合する。ただし××は確認しておらず、一般化には追加測定が必要である。」
 空欄は自分の実測値と根拠で埋める。}
 \noindent 結論（実験後に記入）：\writing{30mm}
 \noindent 改善案・次に確認したい条件：\writing{16mm}
 \section*{参考文献}\label{sec:references}
 \begin{enumerate}
 \item \LabSource
 \item 使用したデータシート・書籍・ウェブ資料（著者、題名、該当箇所、URL、参照日）を追加する。
 \end{enumerate}
}
\newcommand{\prelabclosing}{
 \section*{準備・出典}\label{sec:prelabsource}
 \begin{itemize}[nosep]
 \item 資料、筆記具、記録用紙、電卓と、実験条件・機器設定の記録欄。
 \item 確認問題の解答と、分からない用語・回路点への具体的な質問。
 \end{itemize}
 {\small\noindent\LabSource\par}
}
\newcommand{\simpleflow}[1]{
 \begin{center}
 \begin{tikzpicture}[node distance=5mm and 5mm,
 every node/.style={font=\small},box/.style={draw=labblue,fill=labblue!4,rounded corners=1mm,
 align=center,minimum height=10mm,text width=30mm}]
 #1
 \end{tikzpicture}
 \end{center}
}

\newcommand{\LabCode}{C}
\newcommand{\LabTitle}{電源回路}
\newcommand{\SourcePdf}{IE3_expC_A4.pdf}
\newcommand{\ConclusionGuide}{同じ負荷で比較した直流出力とリプルを根拠に、四条件の違いを要約する。
「全波の方が小さい」だけでなく、どの平滑方式・何mAで・どのリプル定義を使ったかを添える。
負荷特性と充電電流から装置の利点・制約を述べ、近似式の適用範囲を残す。\par}
\begin{document}
\reporttitle
\documentmap{
\tocrow{sec:auto1}{1 目的}{整流・平滑と負荷への供給を実験で確かめる。}
\tocrow{sec:auto2}{2 原理}{整流・平滑と負荷への供給の仕組みと成立条件を整理する。}
\tocrow{sec:auto3}{3 装置と条件}{機器・定数・測定点を確認し、資料の順序で操作する。}
\tocrow{sec:auto4}{4 方法}{機器・定数・測定点を確認し、資料の順序で操作する。}
\tocrow{sec:auto5}{5 結果}{四条件の出力、リプル、充電電流を表や図に記録する。}
\tocrow{sec:auto6}{6 考察：結果を根拠に説明する}{四条件の出力、リプル、充電電流を根拠に、比較と原因の検討を進める。}
\tocrow{sec:closing}{7 結論のまとめ方}{主要な結果、成立範囲、未確認の条件を短くまとめる。}
\tocrow{sec:references}{参考文献}{使用した資料と外部の根拠を示す。}
}

\section{目的}\label{sec:auto1}
半波・全波整流およびコンデンサ入力形・チョーク入力形の平滑回路を比較し、負荷電流による直流出力、リプル、充電電流の変化を調べる。
\section{原理}\label{sec:auto2}
ダイオードで交流を一方向の電流にし、平滑回路で脈動を小さくする。
半波のリプル周波数は電源周波数$f$、全波では$2f$となる。
コンデンサ入力形でリプルが小さいとき、
\[
 \Delta V_{\mathrm{pp}}\simeq\frac{I_L}{f_rC}
\]
と近似できる。資料でいうリプル電圧は波形の最大値と最小値の差の半分であり、$\Delta V_{\mathrm{pp}}/2$として整理する。
チョーク入力形では直列インダクタが電流変動を抑えるが、負荷条件により特性が変わる。
\simpleflow{
\node[box] (a) {交流・変圧器};
\node[box,right=of a] (b) {半波／全波\\整流};
\node[box,right=of b] (c) {C入力／L入力\\平滑};
\node[box,right=of c] (d) {負荷};
\draw[-{Stealth}] (a)--(b);\draw[-{Stealth}] (b)--(c);\draw[-{Stealth}] (c)--(d);
}
\section{装置と条件}\label{sec:auto3}
整流回路、平滑素子、負荷抵抗、直流電圧計、電流計、オシロスコープを使用する。
素子定数、機器番号、プローブ倍率、電源周波数を記録する。
\writing{15mm}
\note{商用電源を含む装置である。電源を切って配線し、装置の指定された測定端子を使う。
オシロスコープの接地端子を任意の回路点につながない。}
\clearpage
\section{方法}\label{sec:auto4}
\begin{enumerate}
\item 資料C-4〜C-9の指定回路を用い、半波・全波とC入力・チョーク入力の四条件を構成する。使用した回路を図示する。
\item 負荷電流を$200,400,600,800\,\mathrm{mA},1\,\mathrm{A}$に合わせ、直流出力電圧を測定する。
\item オシロスコープをAC結合としてリプル波形を読み、ピーク間値とその半分を記録する。
\item $500\,\mathrm{mA}$と$1\,\mathrm{A}$で、整流・平滑に関係する波形とコンデンサ充電電流を記録する。
\item 充電電流測定用$0.1\,\Omega$抵抗の電圧$v_s$から$i_C=v_s/(0.1\,\Omega)$を求める。CH1の入力波形をトリガに用い、CH2で観測点を変える。反転表示の有無も記す。
\end{enumerate}
\section{結果}\label{sec:auto5}
表1：半波整流。リプル欄には半ピーク間値を記入する。
\par\noindent\begin{tabularx}{\linewidth}{YYYYY}
\toprule $I_L$/mA & C入力$V_{\rm DC}$/V & Cリプル/V & L入力$V_{\rm DC}$/V & Lリプル/V\\\midrule
200&&&&\\400&&&&\\600&&&&\\800&&&&\\1000&&&&\\\bottomrule
\end{tabularx}
表2：全波整流。
\par\noindent\begin{tabularx}{\linewidth}{YYYYY}
\toprule $I_L$/mA & C入力$V_{\rm DC}$/V & Cリプル/V & L入力$V_{\rm DC}$/V & Lリプル/V\\\midrule
200&&&&\\400&&&&\\600&&&&\\800&&&&\\1000&&&&\\\bottomrule
\end{tabularx}
\clearpage
図1：$500\,\mathrm{mA}$時の四条件の波形（別紙を追加可）。縦軸・横軸、測定点、反転の有無を示す。
\writing{38mm}
図2：$1\,\mathrm{A}$時の波形および充電電流。
\writing{38mm}
図3：負荷電流と直流出力電圧／リプル電圧の比較グラフ。
\writing{35mm}
\clearpage\section{考察：結果を根拠に説明する}\label{sec:auto6}
\guidepoint{半波と全波でリプルが違う理由}
{同じ負荷・同じ平滑方式で、リプルの周波数と大きさが変わるか。}
{表1・2の同じ負荷行、整流・出力波形、電源周波数。}
{山の回数と次の充電までの時間を比べる。C入力では放電時間の違いを近似式と対応させる。半ピーク間値で比較したことを明記する。}
\guidepoint{負荷増大による出力低下}
{負荷電流増加で直流出力とリプルがどの方向に変わるか。}
{200〜1000 mAの四系列グラフ、各素子値、入力電源条件。}
{一回路の傾向と四回路の比較を分けて記す。電荷の放電、素子・配線の電圧降下を候補として説明し、測定していない要因は確定しない。}
\guidepoint{C入力とチョーク入力の供給の違い}
{直流出力の高さ、リプル、電流の連続性にどのような違いがあるか。}
{同じ整流方式の二平滑条件、500 mAと1 Aの充電電流・出力波形。}
{電圧の山だけでなく電流の流れる区間を示す。チョークの軽負荷条件や損失を考慮し、装置の条件でどちらが適していたか述べる。}
\guidepoint{電流と電圧の時間的な対応}
{出力電圧が回復する区間に充電電流が流れているか。}
{図1・2の共通トリガによる波形、シャント値、CH2反転の有無。}
{導通・放電区間を同じ時間軸で対応させる。シャント電圧をAに換算し、極性と倍率を確認する。瞬時ピークと平均負荷を同じ量として比較しない。}
\subsection{資料の課題・追加検討}
\begin{enumerate}
\item 半波と全波のリプル周波数・大きさを、充放電の時間間隔から説明する。
\item C入力とチョーク入力の負荷特性を比較し、負荷増大時の出力変化を検討する。
\item 出力電圧が最大になる条件と、ダイオードに大きな瞬時電流が流れる条件を波形で説明する。
\item 充電電流波形と平滑コンデンサの電圧波形の時間的な対応を述べる。近似式の適用範囲も検討する。
\end{enumerate}
\writing{22mm}
\reportclosing{四条件の出力電圧とリプルを比較し、適する負荷条件を述べる。}
\end{document}
